\section{Rank deficiencies modulo $p$}\label{sec:content}

In this section we prove Theorem~\ref{thm:content}. The principle is that of the arithmetic method of Hata
and Rhin--Viola, in the form that suits a determinant: if the reduction modulo $p$ of $p^2\cA^0_N$ has
rank at most $n-e$, then by Lemma~\ref{lem:rank-valuation} and Proposition~\ref{prop:valuation} the
determinant $\Delta_N(z)$ gains a factor $p^e$ at every $z$ with denominator prime to $p$, in
particular at every integer. The gain is in the denominator, and by \eqref{eq:QN} it is passed on to
the coefficients of $Q_N$. We therefore need lower bounds for the dimension of the left kernel of
$\cA^0_{N,p}$, that is, explicit row combinations $c\in\F_p^n$ with $c^{\mathsf T}\cA^0_{N,p}=0$. The
tools are the same throughout: the identities for Chebyshev polynomials in characteristic $p$
\begin{equation}\label{eq:cheb-p}
 \begin{gathered}
 T_p=X^p,\qquad U_{p-1}=(X^2-1)^{(p-1)/2},\\
 T_{kp}=T_k(X^p),\qquad U_{kp-1}=(X^2-1)^{(p-1)/2}\,U_{k-1}(X^p)
 \end{gathered}
\end{equation}
in $\F_p[X]$ (the last two follow from $T_{kp}=T_k\circ T_p$ and $U_{kp-1}=(U_{k-1}\circ T_p)\,U_{p-1}$), the relation
$(1-X)^p=1-X^p$, and the digit formula \eqref{eq:filtered-digits}.

\subsection{Setting}
Fix a design, an integer $N\ge1$ and an odd prime $p$ with $H<p^2$. For $0\le r<n$ put $m_r=r-g$, so that
$-g\le m_r\le n-1-g$, and $n-1-g\ge g$. In the notation of \eqref{eq:rows},
\[
 P_r=(1-X)^h\operatorname{rev}_{C-1}T_{m_r},\qquad D_r=(1-X)^h\operatorname{rev}_{C-1}U_{m_r-1},
\]
where $T_{-m}=T_m$ and $U_{-m-1}=-U_{m-1}$, so that $D_g=0$. For $c\in\F_p^n$ write
$S_c=\sum_rc_r\operatorname{rev}_{C-1}T_{m_r}$ and $U_c=\sum_rc_r\operatorname{rev}_{C-1}U_{m_r-1}$, so that $P_c=(1-X)^hS_c$
and $D_c=(1-X)^hU_c$. The lowest degree of $\operatorname{rev}_{C-1}T_m$ is at least $C-1-|m|\ge a:=A=b+q+2g$, and that of
$\operatorname{rev}_{C-1}U_{m-1}$ is at least $a+1$. Put
\[
 Q_c=(X-1)^q\,XWP_c,\qquad \widehat D_c=(X-1)^qD_c,\qquad W=(1-X^2)^{(p-1)/2},
\]
so that, with $\Lambda=g+q+h$, $Q_c$ has its coefficients in $[a+1,\,L+\Lambda+p-1]$ and $\widehat D_c$ in
$[a+1,\,L+\Lambda-1]$. Recall that a vector $c$ lies in the left kernel of $\cA^0_{N,p}$ as soon as the
right side of \eqref{eq:filtered-digits} vanishes for every filtered column $k$ (for the columns with a
carry among the $q+1$ raw columns one splits the sum over the raw columns at the carry; we shall not need
this case).

\subsection{Antisymmetric combinations of rows}
\begin{theorem}\label{thm:T}
Let $2g<p$, $H\le4p$ and $D_T=\min(g,\,n-1-g-2p,\,C-3p)\ge1$. Then
$\nu_N(p)\ge D_T-(n+q+g+h-3p)_+$.
\end{theorem}

The rows used are the differences of rows placed symmetrically about the three centres $g$, $g+p$ and
$g+2p$.

\begin{proof}
For $1\le d\le D_T$ and $k=0,1,2$ let $\delta_k(d)=e_{g+kp+d}-e_{g+kp-d}\in\F_p^n$. These rows exist because
$d\le g$ and $g+2p+d\le n-1$, and the $3D_T$ vectors $\delta_k(d)$ are linearly independent, since their supports
are disjoint: the supports lie in the intervals $[kp,kp+2g]$ of row indices, which are disjoint because
$2g<p$. The row $g+kp\pm d$ has $m=kp\pm d$. By $T_{a+b}-T_{a-b}=2(X^2-1)U_{a-1}U_{b-1}$ and
$U_{a+b-1}-U_{a-b-1}=2T_aU_{b-1}$ together with \eqref{eq:cheb-p},
\begin{align*}
 T_{kp+d}-T_{kp-d}&=2(X^2-1)^{(p+1)/2}\,U_{k-1}(X^p)\,U_{d-1},\\
 U_{kp+d-1}-U_{kp-d-1}&=2\,T_k(X^p)\,U_{d-1}.
\end{align*}
Put $R_d=2(1-X)^hX^{C-d}\operatorname{rev}_{d-1}U_{d-1}$, the value of $D_{\delta_0(d)}$. Reversing at degree $C-1$,
and using that $\operatorname{rev}_{kp}(T_k(X^p))=(\operatorname{rev}_kT_k)(X^p)$, one finds, with $T_0=1$, $T_1=X$ and
$T_2=2X^2-1$,
\[
 D_{\delta_1(d)}=X^{-p}R_d,\qquad D_{\delta_2(d)}=2X^{-2p}R_d-R_d,
\]
\begin{gather*}
 XWP_{\delta_0(d)}=0,\qquad XWP_{\delta_1(d)}=X^{-p}R_d-X^pR_d,\\
 XWP_{\delta_2(d)}=2X^{-2p}R_d-2R_d .
\end{gather*}
The polynomial $R_d$ is supported in degrees $[C-d,H-1]$, which lies in $[3p,4p)$ because $d\le C-3p$ and
$H\le4p$. So at each residue $\ell$ it has a single base-$p$ digit $\varrho_\ell(d)$, in position $u=3$.
For $x=(x_0,x_1,x_2)\in\F_p^3$ the response of the combination $\sum_kx_k\delta_k(d)$ at the raw column
$Jp+\ell$ is therefore, by Lemma~\ref{lem:digits}, $\varrho_\ell(d)\,f_J(x)$ with
\[
 f_J(x)=x_0\beta_{3,J}+x_1(\alpha_{2,J}-\alpha_{4,J}+\beta_{2,J})+x_2(2\alpha_{1,J}-2\alpha_{3,J}+2\beta_{1,J}-\beta_{3,J}),
\]
which depends neither on $d$ nor on $\ell$. The raw columns with $J\ge3$ do not meet the support, since
$3p+\ell\ge C-d\ge C-g=L$. Choose $x^*\ne0$ with $f_1(x^*)=f_2(x^*)=0$; this is two linear equations in three
unknowns. For $c=\sum_ds_d\sum_kx^*_k\delta_k(d)$ the responses at the columns with $J=1,2$ vanish, and the
columns with $J=0$, which have $\ell\ge b$, respond by $f_0(x^*)\sum_ds_d\varrho_\ell(d)$, the coefficient
of $X^{3p+\ell}$ in $\sum_ds_dR_d$. These are at most $(H-3p-b)_+=(n+q+g+h-3p)_+$ linear conditions on
$s=(s_d)$. The map $s\mapsto c$ is injective, so the solutions form a space of dimension at least
$D_T-(n+q+g+h-3p)_+$ in the left kernel.
\end{proof}

\subsection{A reflection of Chebyshev polynomials}
The first ingredient of the next theorem is a reflection formula modulo $p$.

\begin{lemma}\label{lem:u3-reflect}
Let $p$ be an odd prime and $\nu=(p+1)/2$. In $\F_p[X]$, for every integer $m$,
\[
 T_{p-m}=X^pT_m-(X^2-1)^{\nu}\,U_{m-1}.
\]
\end{lemma}

\begin{proof}
Both sides satisfy the three-term recurrence in $m$, so it suffices to check $m=0$ and $m=1$. This
needs \eqref{eq:cheb-p}. For $m=0$ it is $T_p=X^p$. For $m=1$ use $T_{p+1}=XT_p+(X^2-1)U_{p-1}$ and
$T_{p-1}=2XT_p-T_{p+1}$.
\end{proof}

\begin{theorem}\label{thm:U}
Let $2h\le p+1$, $L\le3p\le H$, $n+q+g+h\le3p$ and $2p\le n+q$, and put $n^*=\min(n,2p+2g+1)$. Then
$\nu_N(p)\ge n^*-g-2p-1$.
\end{theorem}

\begin{proof}
We restrict $c$ to the rows $r<n^*$, so that $m_r\le n^*-1-g\le2p+g$. Applying Lemma~\ref{lem:u3-reflect}
with $m=m_r$ and reflecting at degree $p+C-1$ gives, for every such row,
\[
 (1-X^2)^\nu\operatorname{rev}_{C-1}U_{m-1}-X\operatorname{rev}_{C-1}T_m=-X^{p+1}\operatorname{rev}_{C-1}T_{p-m},
\]
and the right side is divisible by $X^{p+C-|p-m|}$, hence by $X^L$ since $|p-m|\le p+g=p+C-L$. Multiplying
by $(1-X)^h$ and summing, for every $c$ supported on these rows
\begin{equation}\label{eq:identityK}
 X^L\mid(1-X^2)^\nu D_c-XP_c .
\end{equation}

Let $V$ be the space of such $c$ with $(1-X^2)^\nu\mid P_c$. Since $(1-X)^h$ divides every $P_r$, this is
$\nu-h$ conditions at $X=1$ and $\nu$ conditions at $X=-1$, that is, at most $p+1-h$ linear conditions.
For $c\in V$ write $P_c=(1-X^2)^\nu Y$ and $\widetilde Q=XY$. Since $(1-X^2)^p=1-X^{2p}$,
\[
 XWP_c=(1-X^{2p})\,\widetilde Q .
\]
The polynomial $P_c$ has no coefficients below $X^{A^*}$, $A^*=C-n^*+g$, and none from $X^H$ on, so
$\widetilde Q$ is supported in $[A^*+1,\,H-p-1]$; and dividing \eqref{eq:identityK} by the unit
$(1-X^2)^\nu$ modulo $X^L$ gives $D_c[k]=\widetilde Q[k]$ for $k<L$.

Now let $H-3p\le\ell<L-2p$. The inequality $L-2p\le A^*+1$ is equivalent to $n^*\le2p+2g+1$, and
$L-2p\le p$ since $L\le3p$. By the supports, $\widetilde Q[\ell]=0$, $\widetilde Q[2p+\ell]=0$ and
$\widetilde Q[3p+\ell]=0$, and $D_c[3p+\ell]=0$ because $3p+\ell\ge H$. So with
$\varkappa=\widetilde Q[p+\ell]$ the base-$p$ digits of $XWP_c$ at the residue $\ell$ are
$(0,\varkappa,0,-\varkappa,0,\dots)$ and those of $D_c$ are $(0,\varkappa,0,0,\dots)$. Hence, by
Lemma~\ref{lem:digits}, the response at every raw column $Jp+\ell$ with such an $\ell$ is $\varkappa$ times
a constant depending on $J$ only: on $V$ these columns, which have $J\in\{0,1,2\}$, span a space of
dimension at most one for each residue $\ell$.

Put $\ell_0=H-3p$ and $\ell_1=L-2p$; the hypotheses give $\ell_0\le b\le\ell_1$. For $\ell_0\le\ell<b$ the raw
columns $j\in[b,L)$ with this residue have $J\in\{1,2\}$, and for $b\le\ell<\ell_1$ they have
$J\in\{0,1,2\}$. So the saving is $2\ell_1-\ell_0-b$, and the number of linear conditions needed to make all
responses on the $n+q=L-b$ raw columns vanish on $V$ is
$(p+1-h)+(n+q)-(2\ell_1-\ell_0-b)$. Substituting $\ell_1=L-2p$, $\ell_0=H-3p$ this is $2p+1+g$. The space of
solutions in the span of $n^*$ rows has dimension at least $n^*-g-2p-1$.
\end{proof}

\subsection{A block of columns: the digits $1$ to $5$}
\begin{theorem}\label{thm:U6}
Let $e=(p+1)/2-q-h\ge0$, $n\le4p+2g+1$, $L\le6p$, and put $\mu_0=\max(q,H+q-6p)$, $\mu_1=L-5p$. Assume
$b+q\le p+\mu_0$. Then $\nu_N(p)\ge\mu_1-\mu_0-e$, that is
\[
 \nu_N(p)\ge\min\Bigl(\frac{p-1}2-g,\ n+b+q+h-\frac{11p+1}2\Bigr).
\]
The first term is the minimum if $6p\le H$ and the second if $6p\ge H$.
\end{theorem}

Unlike Theorems~\ref{thm:T} and~\ref{thm:U}, where the conditions are imposed on the rows, the argument
concerns a block of columns. We first need an identity.

\begin{lemma}\label{lem:K6}
In $\F_p[X]$, for every integer $k$, $2X^p(X^2-1)^{(p-1)/2}T_k-(2X^{2p}-1)U_{k-1}=U_{2p-k-1}$. If
$n\le4p+2g+1$, then for every $c$, $X^L\mid 2XWS_c-(2-X^{2p})U_c$, and consequently
\begin{equation}\label{eq:K6}
 2\,Q_c[i]=2\,\widehat D_c[i]-\widehat D_c[i-2p]\qquad(i<L).
\end{equation}
\end{lemma}

\begin{proof}
For all integers $m,k$ one has $U_{m-1}T_k-T_mU_{k-1}=U_{m-k-1}$, since both sides satisfy the same
recurrence in $k$ and agree for $k=0,1$. Take $m=2p$ and use \eqref{eq:cheb-p}:
$T_{2p}=2X^{2p}-1$ and $U_{2p-1}=2T_pU_{p-1}=2X^p(X^2-1)^{(p-1)/2}$. Reflecting at degree $2p+C-1$, the first
term becomes $2XW\operatorname{rev}_{C-1}T_k$, the second $(2-X^{2p})\operatorname{rev}_{C-1}U_{k-1}$, and the
right side $X^{2p+C-1}U_{2p-k-1}(1/X)$. For $k=m_r\ge-g$ this is divisible by $X^{C+k}$, hence by $X^L$, if
$2p-k-1\ge0$; and if $2p-k-1<0$ it has degree $k-2p-1$ and is divisible by $X^{4p+C-k}$, hence by $X^L$ when
$k\le4p+g$, which holds because $k\le n-1-g\le4p+g$. Summing with weights $c_r$ and multiplying by
$(X-1)^q(1-X)^h$ gives \eqref{eq:K6}.
\end{proof}

\begin{proof}[Proof of Theorem~\ref{thm:U6}]
Let $M=[\mu_0,\mu_1)$ and $w=\mu_1-\mu_0$; we may assume $w>e$. The block consists of the filtered columns
$k(J,\mu)=Jp+\mu-(b+q)$ for $1\le J\le5$ and $\mu\in M$. These indices are non-negative by the hypothesis
$b+q\le p+\mu_0$, and smaller than $n$ because $5p+\mu_1-1=L-1$; so there are $5w$ of them. As $\mu<p$, formula
\eqref{eq:filtered-digits} applies: the response of $c$ at $k(J,\mu)$ is
$\sum_u\alpha_{u,J}Q_c[up+\mu]+\beta_{u,J}\widehat D_c[up+\mu]$.

\emph{Supports.} For $\mu\in M$ we have $p+\mu\le p+\mu_1-1=L-4p-1\le a$ by $n\le4p+2g+1$, so the digits
$u=0,1$ of $Q_c$ and $\widehat D_c$ vanish. From $\mu_0\ge H+q-6p=L+\Lambda-6p$ we get that the digits $u\ge6$ of
$\widehat D_c$ and $u\ge7$ of $Q_c$ vanish. For $u\le5$ we have $up+\mu\le L-1$, so \eqref{eq:K6} applies.

\emph{From $Q$ to $\widehat D$.} For $u=2,\dots,5$, \eqref{eq:K6} gives
$\widehat D_c[up+\mu]=Q_c[up+\mu]+\tfrac12\widehat D_c[(u-2)p+\mu]$, and for $u=2,3$ the last index is at most $a$,
so the term vanishes. Hence if $Q_c[up+\mu]=0$ for $u=2,\dots,5$ and all $\mu\in M$, then by induction
$\widehat D_c[up+\mu]=0$ for $u=2,\dots,5$, and all digits of $\widehat D_c$ in $M$ vanish.

\emph{A relation from the Frobenius map.} We have
$Q_c=(-1)^qX(1-X)^{q+h+(p-1)/2}(1+X)^{(p-1)/2}S_c$, and $q+h+(p-1)/2+e=p$. Since $(1-X)^p=1-X^p$,
\[
 (1-X)^e\,Q_c=(1-X^p)\,Y_c,\qquad Y_c=(-1)^qX(1+X)^{(p-1)/2}S_c,
\]
and $Y_c$ is supported in $[a+1,\,C+(p-1)/2]$. Comparing the coefficients at $up+\mu$, subtracting the
coefficient at $(u-1)p+\mu$, and summing over $u=2,\dots,6$, the right side telescopes to
$Y_c[6p+\mu]-Y_c[p+\mu]$, which vanishes for $\mu\ge\mu_0+e$: the second term because $p+\mu\le a$, the
first because $6p+\mu\ge H+q+e=C+(p+1)/2>\deg Y_c$. So, if $Q_c[up+\mu']=0$ for $u=2,\dots,5$ and all
$\mu'\in M$, then
\[
 \sum_{t=0}^e(-1)^t\binom et\,Q_c[6p+\mu-t]=0\qquad(\mu_0+e\le\mu<\mu_1).
\]
This is a recurrence with unit leading coefficient: if $Q_c[6p+\mu]=0$ for the first $e$ values
$\mu=\mu_0,\dots,\mu_0+e-1$, then it vanishes for all $\mu\in M$.

\emph{Counting.} Impose on $c$ that the $n-5w$ filtered columns outside the block vanish, that
$Q_c[up+\mu]=0$ for $2\le u\le5$ and $\mu\in M$ ($4w$ conditions), and that $Q_c[6p+\mu]=0$ for the first $e$
values of $\mu$. By the three steps above all the responses of the block columns vanish. These are at
most $n-5w+4w+e=n-(w-e)$ linear conditions, so the left kernel has dimension at least $w-e$. Finally,
$w-e$ equals $\mu_1-q-e=n+b+q+h-(11p+1)/2$ if $\mu_0=q$ and $(p-1)/2-g$ if $\mu_0=H+q-6p$.
\end{proof}

\subsection{The second window}
\begin{theorem}\label{thm:S}
Let $e=(p+1)/2-q-h\ge0$ and $n\le4p+2g+1$, and put
\[
 m_0=\max(b+q,\,L+\Lambda-4p),\qquad m_1=\min(p,\,L-3p,\,a+1),\qquad a=b+q+2g .
\]
Then $\nu_N(p)\ge m_1-m_0-e$.
\end{theorem}

\begin{proof}
Let $M=[m_0,m_1)$ and $w=m_1-m_0>e$. For $\mu\in M$ we have $\mu\le a$, so the digit $u=0$ of $Q_c$ and
$\widehat D_c$ vanishes; $\mu\ge L+\Lambda-4p$ gives that the digits $u\ge4$ of $\widehat D_c$ and $u\ge5$ of
$Q_c$ vanish; and $3p+\mu<L$, so that \eqref{eq:K6} gives, with $D_i=\widehat D_c[ip+\mu]$ and
$Q_i=Q_c[ip+\mu]$,
\[
 Q_1=D_1,\qquad Q_2=D_2,\qquad Q_3=D_3-\tfrac12D_1 .
\]
So every response of the filtered matrix at a column with residue $\mu\in M$ is a combination of the four
functionals $D_1(\mu),D_2(\mu),D_3(\mu),Q_4(\mu)$ of $c$. For $\mu\in M$ the columns $k=\mu-(b+q)$,
$p+\mu-(b+q)$, $2p+\mu-(b+q)$ and $3p+\mu-(b+q)$ exist, since $b+q\le\mu<p$ and $3p+\mu<L$. So there are at least
$4w$ block columns, and their rows lie in the span $F$ of the $4w$ functionals.

We find $w-e$ linear relations among these functionals. As in the proof of Theorem~\ref{thm:U6},
$(1-X)^eQ_c=(1-X^p)Y_c$ with $Y_c$ a polynomial. Put $\sigma(\ell)=\sum_{u\ge0}Q_c[up+\ell]$ and sum the
coefficients of this identity over the residue class of $\mu\ge e$ modulo $p$. The right side gives
$\sum_u\bigl(Y_c[up+\mu]-Y_c[(u-1)p+\mu]\bigr)=0$, since $Y_c$ is a polynomial, and the left side gives
$\sum_{t=0}^e(-1)^t\binom et\sigma(\mu-t)$. For $\mu\in M$ with $\mu-e\ge m_0$ all $\mu-t$ lie in $M$, and there
$\sigma=Q_1+Q_2+Q_3+Q_4=D_1+D_2+D_3-\tfrac12D_1+Q_4$. So
\[
 \sum_{t=0}^e(-1)^t\binom et\bigl(\tfrac12D_1+D_2+D_3+Q_4\bigr)(\mu-t)=0\qquad(m_0+e\le\mu<m_1),
\]
and these $w-e$ relations are independent, since $Q_4(\mu)$ occurs with coefficient $1$ in the one for $\mu$
and only with lower arguments in the others. Hence $\dim F\le4w-(w-e)$, and the rank of the matrix is at most
$(n-4w)+3w+e=n-(w-e)$.
\end{proof}

The proof does not use a determinant or a lower bound for $N$. Unlike the proof of Theorem~\ref{thm:U6}, the
telescoping here is over all digits, so no vanishing of $Y_c$ at the ends is needed.

\subsection{The function $\kappa$ for the design $D$}
For $D$, $L=70N$, $H=78N$, $a=22N$, $\Lambda=11N$ and $q+h=6N$. Writing $x=p/N$, one checks the hypotheses of
the theorems on the intervals of the table of Section~\ref{sec:overview}. The hypotheses
$e\ge0$ and $n\le4p+2g+1$ both read $x\ge12$, and $L\le6p$ reads $x\ge35/3$; so Theorems~\ref{thm:U6}
and~\ref{thm:S} apply for $x\ge12$. For $12\le x\le13$ one has $\mu_0=H+q-6p=81N-6p$, so the first branch
of Theorem~\ref{thm:U6} applies, and for $x\ge13$ one has $\mu_0=3N$ and the second branch, which is positive
until $x=146/11$. In Theorem~\ref{thm:S}, $m_1=p$ for $x\le35/2$ and $m_1=70N-3p$ above, and
$m_0=81N-4p$ for $x\le69/4$ and $12N$ above; the value $m_1-m_0-e$ is $9x/2-75$, $x/2-6$ and
$64-7x/2$ on the three intervals $(50/3,69/4]$, $(69/4,35/2]$, $(35/2,128/7]$, and it is non-positive
outside $(50/3,128/7)$. Theorem~\ref{thm:T} applies for $x>10$ and $x\ge39/2$, and
$D_T=\min(5N,52N-1-2p,75N-3p)$; it gives $3p-(n+q+h)=3x-64$ up to $x=23$ and the cap $g=5N$ up to
$x=70/3$. From $x=70/3$ the hypotheses of Theorem~\ref{thm:U} hold up to $x=26$ (they read
$3p\ge L$, $3p\le H$, $2p\le n+q$ and $2h\le p+1$), and give $\min(n,2p+2g+1)-g-2p-1$, which is $5N$ for
$x\le24$ and $53N-2p-1$ above. In each case the value of the theorem is the value in the table minus at
most one unit, from the integer rounding. At a prime the largest of the values that apply is used.
The area of $\kappa$ is the sum of the areas of the eight cells, which are listed in
Section~\ref{sec:overview}.

\begin{remark}
Numerical computations of the exact rank of $\cA^0_{N,p}$ for several designs and primes show that the
bounds of Theorem~\ref{thm:content} are attained in the windows where they were tested, and that other
families add kernel where the windows overlap. These computations are not part of the proof. An upper bound
for the nullity is not proved, and none is needed.
\end{remark}
